\chapter{角动量与自旋}
\label{chap:angular-momentum-spin}
\index{角动量}
\index{自旋}

\section{角动量在量子力学中的作用}
\label{sec:angular-momentum-role}

在经典力学中,角动量可以看作旋转的 Hamilton 生成元(命题~\ref{pro:angular-momentum-flow}).当系统具有旋转对称性时,角动量尤其有用,因为它是守恒量(推论~\ref{cor:poisson-conserved-quantity}).在量子力学中,角动量仍然是旋转的“生成元”:按照 Stone 定理(定理~\ref{thm:stone}),它是一参数酉旋转算子群的无穷小生成元.具有旋转对称性的系统中,量子角动量同样守恒;也就是说,若 Hamilton 算子 $\widehat H$ 在旋转下不变,则它与角动量算子对易,角动量算子因而是量子力学意义下的运动常量.

$\mathbb R^3$ 中粒子的经典角动量向量的各分量,在 Poisson 括号下满足简单的交换关系(第 2 章习题 19).我们将看到,这些正是旋转群 $SO(3)$ 的 Lie 代数 $\mathfrak{so}(3)$ 的交换关系.若 $\widehat H$ 与角动量的每个分量对易,则 $\widehat H$ 的每个特征子空间(即给定 $\lambda$ 时方程 $\widehat H\psi=\lambda\psi$ 的解空间)在角动量算子下不变.因此,该特征子空间构成 $\mathfrak{so}(3)$ 的表示.通过分类 $\mathfrak{so}(3)$ 的不可约有限维表示,可以得到旋转不变情形下解空间的大量信息.

具体而言,$\mathfrak{so}(3)$ 的表示论允许我们完全确定解 $\psi(x)$ 的角向依赖,只留下待求的径向依赖.这把独立变量的个数从三个($\mathbb R^3$ 中的三个坐标)降为一个(极坐标中的半径 $r$),从而把问题化为常微分方程.

从 Lie 代数表示的角度理解角动量,也为理解自旋作准备.带自旋的 $\mathbb R^3$ 中粒子的 Hilbert 空间是 $L^2(\mathbb R^3)$ 与有限维向量空间 $V$ 的张量积,其中 $V$ 承载旋转群 $SO(3)$ 的不可约作用.在这里,适当的“作用”概念是 $SO(3)$ 的射影表示,即满足 $SO(3)$ 关系但允许相差相位因子的算子族.允许这些相位因子,是因为物理上只相差一个常数的两个向量表示同一物理状态.由命题~\ref{prop:projective-representation-lie-lift},$SO(3)$ 的每个射影表示都能在 Lie 代数 $\mathfrak{so}(3)$ 层面去射影化;反过来,$\mathfrak{so}(3)$ 的每个不可约普通表示都给出其万能覆盖 $SU(2)$ 的表示,继而给出 $SO(3)$ 的射影表示.因此,空间 $V$ 的可能性与 $\mathfrak{so}(3)$ 的不可约表示一一对应.对于“半整数自旋”,$V$ 不承载 $SO(3)$ 的普通表示.

\section{$\mathbb R^3$ 中的角动量算子}
\label{sec:angular-momentum-operators}
\label{sec:angular-momentum}

回忆第~\ref{sec:classical-angular-momentum} 节,$\mathbb R^3$ 中粒子的经典角动量为 $J=x\times p$,例如 $J_3=x_1p_2-x_2p_1$.如第~\ref{sec:quantum-axioms} 节,我们引入量子对应物——其分量为算子的“向量” $\widehat J$:
\[
\widehat J=X\times P.
\]
例如,$\widehat J_1=X_2P_3-X_3P_2$.每个角动量分量只涉及位置与动量的不同分量,而这些分量彼此对易.因此在 $\widehat J_3$ 中,写成 $X_2P_3$ 或 $P_3X_2$ 没有区别.

角动量算子是无界算子,只定义在 $L^2(\mathbb R^3)$ 的稠密子空间上.暂时不指定它们的定义域;下一节会处理这一点.不过可以看到,每个角动量算子的定义域都包含 Schwartz 空间 $\mathcal S(\mathbb R^3)$(定义~\ref{app:def:schwartz}).

由正则交换关系(第 3 章习题 10)可得
\[
\frac1{i\hbar}[\widehat J_1,\widehat J_2]=\widehat J_3,
\qquad
\frac1{i\hbar}[\widehat J_2,\widehat J_3]=\widehat J_1,
\qquad
\frac1{i\hbar}[\widehat J_3,\widehat J_1]=\widehat J_2.
\label{eq:angular-momentum-commutators}
\]
这些关系与经典角动量分量之间的 Poisson 括号关系相符.

将 $\widehat J_3$ 写开,有
\[
(\widehat J_3\psi)(x)=-i\hbar\left(x_1\frac{\partial}{\partial x_2}-x_2\frac{\partial}{\partial x_1}\right)\psi(x)
=-i\hbar\left.\frac{d}{d\theta}\psi(R_\theta x)\right|_{\theta=0},
\label{eq:j3-rotation-derivative}
\]
其中 $R_\theta$ 表示 $(x_1,x_2)$ 平面内逆时针旋转 $\theta$ 的旋转矩阵,$\widehat J_1,\widehat J_2$ 也有类似表达式.这说明角动量算子与旋转密切相关,正如经典角动量与旋转密切相关.

定义无量纲角动量算子
\[
\widetilde J_j=\frac1\hbar\widehat J_j.
\label{eq:dimensionless-angular-momentum}
\]
由于 Planck 常数与角动量具有相同单位,$\widetilde J_j$ 不依赖于单位选择;我们称它们为无量纲角动量算子.

\section{Lie 代数观点下的角动量}
\label{sec:angular-momentum-lie-algebra}

先考察旋转群 $SO(3)$ 在 $L^2(\mathbb R^3)$ 上的自然作用.

\begin{definition}[旋转群的自然表示]\label{def:rotation-representation}
对每个 $R\in SO(3)$,定义 $\Pi(R):L^2(\mathbb R^3)\to L^2(\mathbb R^3)$ 为
\[
(\Pi(R)\psi)(x)=\psi(R^{-1}x).
\label{eq:rotation-representation}
\]
\end{definition}

\begin{proposition}[旋转的酉表示]\label{prop:rotation-unitary-representation}
对每个 $R\in SO(3)$,$\Pi(R)$ 是酉算子;此外,$\Pi:SO(3)\to U(L^2(\mathbb R^3))$ 是强连续同态.
\end{proposition}

\begin{proof}
Lebesgue 测度在旋转下不变,因此每个 $\Pi(R)$ 都酉.直接检查可得 $\Pi(R_1R_2)=\Pi(R_1)\Pi(R_2)$;定义中必须使用 $\psi(R^{-1}x)$ 而不是 $\psi(Rx)$ 才能得到这一关系.强连续性可仿照例~\ref{ex:translation-unitary-group} 的证明验证.
\end{proof}

回忆第 16 章中对 $SO(3)$ Lie 代数 $\mathfrak{so}(3)$ 的计算,以及其中的基 $\{F_1,F_2,F_3\}$.对 $X\in\mathfrak{so}(3)$,令 $\pi(X)$ 为满足
\[
\Pi(e^{tX})=e^{t\pi(X)}
\label{eq:rotation-lie-representation}
\]
的斜自伴算子.

\begin{proposition}[角动量与旋转 Lie 代数]\label{prop:angular-momentum-lie-action}
每个 $\pi(F_j)$ 的定义域包含 Schwartz 空间 $\mathcal S(\mathbb R^3)$,并且在该空间上
\[
\widehat J_j=i\hbar\,\pi(F_j).
\label{eq:angular-momentum-lie-identification}
\]
\end{proposition}

\begin{proof}
以 $\widehat J_3$ 为例,$e^{tF_3}$ 是 $(x_1,x_2)$ 平面中的逆时针旋转.若 $\psi\in\mathcal S(\mathbb R^3)$,式~\eqref{eq:j3-rotation-derivative} 中定义导数的极限在 $L^2$ 意义下成立.结合式~\eqref{eq:rotation-representation} 右侧的逆矩阵,得到 $\widehat J_3=i\hbar\pi(F_3)$.$\widehat J_1,\widehat J_2$ 的计算相同.
\end{proof}

虽然精确确定每个角动量算子的定义域并不容易,但由命题~\ref{prop:infinite-representation-finite-subspace},若 $\psi$ 属于 $L^2(\mathbb R^3)$ 中某个在旋转下不变的有限维子空间,则 $\psi$ 属于每个 $\widehat J_j$ 的定义域.

\section{$\mathfrak{so}(3)$ 的不可约表示}
\label{sec:irreducible-so3-representations}

本节分类 $\mathfrak{so}(3)$ 的有限维不可约表示(按同构分类).所有表示均在复数域上,并且维数至少为 $1$.继续使用第 16 章中的基 $\{F_1,F_2,F_3\}$.

\begin{theorem}[$\mathfrak{so}(3)$ 不可约表示的标准形式]\label{thm:so3-irreducible-standard-form}
设 $\pi:\mathfrak{so}(3)\to\mathfrak{gl}(V)$ 是有限维不可约表示.定义 $V$ 上的算子
\[
L_+=i\pi(F_1)-\pi(F_2),\qquad
L_-=i\pi(F_1)+\pi(F_2),\qquad
L_3=i\pi(F_3).
\]
令 $l=\frac12(\dim V-1)$,于是 $\dim V=2l+1$.则存在 $V$ 的一组基 $v_0,v_1,\ldots,v_{2l}$,使
\[
L_3v_j=(l-j)v_j,
\]
且
\[
L_-v_j=
\begin{cases}v_{j+1},&j<2l,\\0,&j=2l,\end{cases}
\qquad
L_+v_j=
\begin{cases}j(2l+1-j)v_{j-1},&j>0,\\0,&j=0.
\end{cases}
\label{eq:so3-irreducible-action}
\]
\end{theorem}

因此,数 $l$ 完全决定不可约表示的结构.由于 $\dim V$ 是正整数,$l$ 必须取
\[
l=0,\frac12,1,\frac32,\ldots.
\label{eq:so3-spin-values}
\]

\begin{definition}[自旋]\label{def:spin}
若 $(\pi,V)$ 是 $\mathfrak{so}(3)$ 的有限维不可约表示,则算子 $L_3=i\pi(F_3)$ 的最大特征值称为 $(\pi,V)$ 的自旋.等价地,自旋 $l$ 是满足 $\dim V=2l+1$ 的唯一数.
\end{definition}

\begin{theorem}[$\mathfrak{so}(3)$ 表示的存在唯一性]\label{thm:so3-irreducible-existence-uniqueness}
对任意 $l=0,\frac12,1,\frac32,\ldots$,存在一个维数为 $2l+1$ 的 $\mathfrak{so}(3)$ 不可约表示;同一维数的任意两个不可约表示都同构.
\end{theorem}

定理~\ref{thm:so3-irreducible-existence-uniqueness} 只断言每个 $l$ 都存在 $\mathfrak{so}(3)$ 的表示.下一节将看到,$\mathfrak{so}(3)$ 的不可约表示来自 $SO(3)$ 的表示,当且仅当 $l$ 是整数;半整数自旋表示虽然不来自 $SO(3)$ 的普通表示,却在量子物理中发挥重要作用.

由命题~\ref{prop:angular-momentum-lie-action},$L_3$ 可以看作空间 $V$ 上无量纲角动量第三分量的模拟.稍后把定理~\ref{thm:so3-irreducible-standard-form} 应用于 $L^2(\mathbb R^3)$ 中的旋转不变子空间时,$L_3$ 正是 $\widetilde J_3$ 在该子空间上的限制.

定理~\ref{thm:so3-irreducible-standard-form} 与量子谐振子的分析十分相似:都有某个算子的特征向量链,以及升降特征值的升降算子.不同之处在于,谐振子的链从基态开始并向一个方向无限延伸,而 $\mathfrak{so}(3)$ 表示的链在两个方向都有限.链从 $L_3$ 最大特征值的特征向量 $v_0$ 开始,此时升算子 $L_+$ 消灭 $v_0$.证明的关键在于确定链在 $L_3$ 较低特征值方向如何终止,而不违背 $L_3,L_+,L_-$ 的交换关系.

\begin{proof}[定理~\ref{thm:so3-irreducible-standard-form} 的证明]
由于 $\pi$ 是 Lie 代数同态,$\pi(F_j)$ 满足与 $F_j$ 相同的交换关系.由此得到
\[
[L_3,L_+]=L_+,qquad [L_3,L_-]=-L_-,qquad [L_+,L_-]=2L_3.
\label{eq:so3-L-commutators}
\]
由于在代数闭域 $\mathbb C$ 上,$L_3$ 至少有一个特征向量 $v$,设其特征值为 $\lambda$.由第一条交换关系,
\[
L_3L_+v=(L_+L_3+L_+)v=(\lambda+1)L_+v.
\]
所以 $L_+v$ 要么为零,要么是特征值为 $\lambda+1$ 的特征向量.这解释了“升算子”的名称.

反复应用 $L_+$,特征值每次增加 $1$,直到得到零向量;因为 $L_3$ 只有有限多个特征值,存在 $k\geq0$ 使 $(L_+)^kv\neq0$ 而 $(L_+)^{k+1}v=0$.令 $v_0=(L_+)^kv$,并令 $\mu=\lambda+k$,则 $L_+v_0=0$、$L_3v_0=\mu v_0$.定义
\[
v_j=(L_-)^jv_0,qquad j=0,1,2,\ldots.
\]
由第二条交换关系,$L_-$ 使 $L_3$ 的特征值降低 $1$,所以 $L_3v_j=(\mu-j)v_j$.

我们声称对 $j\geq1$,
\[
L_+v_j=j(2\mu+1-j)v_{j-1}.
\label{eq:so3-raising-formula}
\]
这个公式可用第三条交换关系对 $j$ 归纳证明(见习题~\ref{ex:so3-raising-induction}).由于 $L_3$ 只有有限多个特征值,$v_j$ 最终必须为零.取 $N$ 使 $v_N\neq0$、$v_{N+1}=0$.将式~\eqref{eq:so3-raising-formula} 用于 $j=N+1$ 得
\[
0=L_+v_{N+1}=(N+1)(2\mu-N)v_N.
\]
因此 $\mu=N/2$.

令 $l=N/2$,则式~\eqref{eq:so3-irreducible-action} 成立.由于 $v_j$ 是具有不同特征值的特征向量,它们线性无关;它们的张成空间在 $L_+,L_-,L_3$ 下不变,因而在整个 $\mathfrak{so}(3)$ 下不变.不可约性迫使该张成空间等于 $V$,所以这些向量构成基,且 $\dim V=N+1=2l+1$.
\end{proof}

\begin{proof}[定理~\ref{thm:so3-irreducible-existence-uniqueness} 的证明]
给定 $l=0,\frac12,1,\ldots$,取以 $v_0,\ldots,v_{2l}$ 为基的向量空间 $V$,并用式~\eqref{eq:so3-irreducible-action} 定义 $L_+,L_-,L_3$ 的作用.直接计算(见习题~\ref{ex:so3-representation-commutators})可验证这些算子满足式~\eqref{eq:so3-L-commutators},所以它们定义 $\mathfrak{so}(3)$ 的表示.

设 $W\neq\{0\}$ 是不变子空间,取非零 $w=\sum_{j=0}^{2l}a_jv_j\in W$,令 $j_0$ 为 $a_j\neq0$ 的最大指标.根据式~\eqref{eq:so3-irreducible-action},$(L_+)^{j_0}w$ 是 $v_0$ 的非零倍数,所以 $v_0\in W$;反复应用 $L_-$ 得 $v_j\in W$ 对所有 $j$ 成立,故 $W=V$.表示不可约.

定理~\ref{thm:so3-irreducible-standard-form} 说明任意同一维数的不可约表示都有这种标准基;把一个表示的标准基映到另一个表示的标准基就得到同构.
\end{proof}

\begin{proposition}[不可约表示的酉内积]\label{prop:so3-unitary-inner-product}
若 $\pi:\mathfrak{so}(3)\to\mathfrak{gl}(V)$ 是不可约表示,则 $V$ 上存在一个内积,且除乘以正常数外唯一,使对所有 $X\in\mathfrak{so}(3)$,$\pi(X)$ 斜自伴.
\end{proposition}

\begin{proof}
由 $L_3,L_+,L_-$ 的定义,$\pi(X)$ 全部斜自伴等价于 $L_3$ 自伴且 $L_+$ 与 $L_-$ 互为伴随.由于 $v_j$ 是 $L_3$ 的不同特征值特征向量,要使 $L_3$ 自伴,$v_j$ 必须正交;反之,若它们正交,$L_3$ 就自伴.

再要求 $L_+^*=L_-$.由式~\eqref{eq:so3-irreducible-action} 与交换关系计算得到
\[
\langle v_j,v_j\rangle=j(2l-j+1)\langle v_{j-1},v_{j-1}\rangle,
\qquad j=1,\ldots,2l.
\label{eq:so3-inner-product-normalization}
\]
这些系数均为正,所以总能把各 $v_j$ 归一化到满足该条件.除 $v_0$ 的归一化外没有自由度,故内积除乘以正常数外唯一.
\end{proof}

\begin{proposition}[Casimir 算子]\label{prop:so3-casimir}
设 $(\pi,V)$ 是维数 $2l+1$ 的不可约 $\mathfrak{so}(3)$ 表示,定义
\[
C_\pi=\pi(F_1)^2+\pi(F_2)^2+\pi(F_3)^2.
\]
则对所有 $v\in V$,
\[
C_\pi v=-l(l+1)v.
\]
\end{proposition}

\begin{proof}
见习题~\ref{ex:so3-casimir}.
\end{proof}

如果考察定理~\ref{thm:so3-irreducible-standard-form} 的证明,会发现唯一使用不可约性的地方,是为了说明 $v_0,\ldots,v_{2l}$ 的张成空间等于整个 $V$.因此有如下结果.

\begin{proposition}[一般有限维表示中的不可约链]\label{prop:so3-cyclic-chain}
设 $(\pi,V)$ 是任意有限维 $\mathfrak{so}(3)$ 表示,不必不可约.若 $v_0\neq0$ 满足 $L_+v_0=0$、$L_3v_0=\lambda v_0$,则 $\lambda$ 是非负整数或半整数 $l$.此外,令
\[
v_j=(L_-)^jv_0,qquad j=0,1,\ldots,2l,
\]
则这些向量张成 $V$ 的一个维数为 $2l+1$ 的不可约不变子空间,并且 $L_+,L_-,L_3$ 在这些向量上的作用由定理~\ref{thm:so3-irreducible-standard-form} 给出.
\end{proposition}

一般地,给定 Lie 代数表示 $(\pi,V)$ 和非零向量 $v_0$,若包含 $v_0$ 的最小不变子空间是整个 $V$,则称 $v_0$ 是 $V$ 的循环向量.在命题~\ref{prop:so3-cyclic-chain} 中,$v_0$ 确实是 $W=\operatorname{span}(v_0,\ldots,v_{2l})$ 的循环向量;但表示有循环向量一般并不意味着表示不可约.这里 $W$ 的不可约性来自 $v_0$ 所满足的特殊条件.

\section{$SO(3)$ 的不可约表示}
\label{sec:irreducible-so3-group-representations}

我们已经分类了 Lie 代数 $\mathfrak{so}(3)$ 的不可约表示,现在转而分类群 $SO(3)$ 的表示.由于 $SO(3)$ 是连通的(见第 16 章习题 13),命题~\ref{prop:group-lie-invariant-subspaces} 告诉我们:$SO(3)$ 的表示不可约,当且仅当其相应的 Lie 代数表示不可约;两个 $SO(3)$ 表示同构,当且仅当相应的 Lie 代数表示同构.因此,要按同构分类 $SO(3)$ 的不可约表示,只需确定 $\mathfrak{so}(3)$ 的哪些不可约表示来自群 $SO(3)$ 的表示.

\begin{proposition}[$SO(3)$ 的不可约表示]\label{prop:so3-group-representations}
设 $\pi_l:\mathfrak{so}(3)\to\mathfrak{gl}(V)$ 是 $\mathfrak{so}(3)$ 的不可约表示,并令
\[
l:=\frac12(\dim V-1)
\]
为其自旋.若 $l$ 是整数(亦即 $V$ 的维数为奇数),则存在 $SO(3)$ 的表示 $\Pi_l:SO(3)\to\operatorname{GL}(V)$,使得 $\Pi_l$ 与 $\pi_l$ 满足定理~\ref{thm:group-homomorphism-differential} 中的关系.若 $l$ 是半整数(亦即 $V$ 的维数为偶数),则不存在这样的表示 $\Pi_l$.
\end{proposition}

由此结果和命题~\ref{prop:group-lie-invariant-subspaces} 可知,$SO(3)$ 的不可约表示恰好是那些 $l$ 为整数时的 $\Pi_l$.

\begin{proof}
先设 $l$ 为半整数.按照基 $\{v_j\}$,$L_3$ 是对角的,其特征值为半整数.因此
\[
 e^{2\pi\pi_l(F_3)}=e^{-2\pi iL_3}=-I.
\]
(这里 $\pi$ 出现在 $\pi_l$ 前面时表示圆周率 $\pi=3.14\ldots$.)另一方面,根据例~\ref{ex:rotation-matrix-exponential} 的简单变形,$\mathfrak{so}(3)$ 中的矩阵 $F_3$ 满足 $e^{2\pi F_3}=I$.于是,若相应的表示 $\Pi_l$ 存在,则应有
\[
 \Pi_l(I)=\Pi_l(e^{2\pi F_3})=e^{2\pi\pi_l(F_3)}=-I,
\]
矛盾.

再设 $l$ 为整数.利用例~\ref{ex:su2-so3-lie-algebras} 的证明中描述的 $\mathfrak{su}(2)$ 与 $\mathfrak{so}(3)$ 之间的同构 $\varphi$,它把 $\mathfrak{su}(2)$ 的基 $\{E_1,E_2,E_3\}$ 映到 $\mathfrak{so}(3)$ 的基 $\{F_1,F_2,F_3\}$.令
\[
 \pi_l'(X)=\pi_l(\varphi(X)),
\]
得到 $\mathfrak{su}(2)$ 的一个表示.由于 $SU(2)$ 是单连通的,定理~\ref{thm:lie-algebra-homomorphism-integrates} 告诉我们,存在 $SU(2)$ 的表示 $\Pi_l'$,并且它与 $\pi_l'$ 具有通常的关系.现在计算得
\[
 \Pi_l'(-I)=\Pi_l'(e^{2\pi E_1})
 =e^{2\pi\pi_l'(E_1)}
 =e^{2\pi\pi_l(F_1)}
 =e^{-2\pi iL_3}=I,
\]
因为 $L_3$ 的特征值都是整数.

由例~\ref{ex:su2-universal-cover-so3},存在从 $SU(2)$ 到 $SO(3)$ 的满同态 $\Phi$,其相应的 Lie 代数同态为 $\varphi$,并且 $\ker\Phi=\{I,-I\}$.由于 $\Pi_l'$ 的核包含 $\{I,-I\}$,映射 $\Pi_l'$ 可以经由 $SO(3)$ 分解,得到 $SO(3)$ 的表示 $\Pi_l$,满足
\[
 \Pi_l'=\Pi_l\circ\Phi.
\]
根据第 16 章习题 10,相应的 $\mathfrak{so}(3)$ 表示 $\sigma_l$ 满足 $\pi_l'=\sigma_l\circ\varphi$,从而
\[
 \sigma_l=\pi_l'\circ\varphi^{-1}=\pi_l.
\]
所以 $\Pi_l$ 正是所需的 $SO(3)$ 表示.
\end{proof}

\section{在 $L^2(S^2)$ 中实现这些表示}
\label{sec:representations-in-l2-sphere}

本节不采用物理文献中的传统处理方式:我们把“球谐函数”看作 $\mathbb R^3$ 中某些多项式在单位球面上的限制,而不是用球面的角坐标来描述它们.这样做可以避免极坐标中的一些繁琐计算,并且很容易推广到更高维.

回忆第~\ref{sec:angular-momentum-lie-algebra} 节,$SO(3)$ 在 $L^2(\mathbb R^3)$ 上有自然的酉表示 $\Pi$,其作用为
\[
 (\Pi(R)\psi)(x)=\psi(R^{-1}x).
\]
在研究量子氢原子等旋转不变问题时,理解 $L^2(\mathbb R^3)$ 中那些在 $\Pi$ 下不变、且 $\Pi$ 的限制不可约的有限维子空间 $V$ 的结构会很有用.

若用极坐标写出 $\mathbb R^3$ 上的函数,则 $SO(3)$ 只作用于角变量.因此也应考虑 $SO(3)$ 在 $L^2(S^2)$ 上的作用,公式与 $L^2(\mathbb R^3)$ 上相同,即
\[
 (\Pi(R)\psi)(x)=\psi(R^{-1}x),\qquad x\in S^2.
\]
计算 $L^2(S^2)$ 的范数时,我们使用在 $SO(3)$ 作用下不变的球面面积测度.找到 $L^2(S^2)$ 中的不变子空间后,正如下一节所见,很容易构造 $L^2(\mathbb R^3)$ 中相应的不变子空间.

本节关注 $\mathbb R^3$ 上的调和多项式,即满足 $\Delta p=0$ 的多项式,其中 $\Delta$ 是 Laplace 算子.由于我们始终在复数域上考虑表示,允许这些多项式具有复系数.

\begin{definition}[次数为 $l$ 的球谐函数空间]\label{def:spherical-harmonics-space}
设 $l$ 是非负整数.令 $V_l$ 为 $\mathbb R^3$ 上次数为 $l$ 的齐次调和多项式限制到 $S^2$ 后得到的函数空间.称 $V_l$ 为次数为 $l$ 的球谐函数空间.
\end{definition}

注意,若 $p$ 是某个次数 $l$ 的齐次多项式,则它在 $S^2$ 上的限制恒等于零,当且仅当 $p$ 本身恒等于零.事实上,如果 $p$ 是 $l$ 次齐次多项式,并且在 $S^2$ 上为零,那么对所有 $x\neq0$,都有
\[
 p(x)=|x|^l p\left(\frac{x}{|x|}\right)=0,
\]
再由连续性可知 $p$ 在所有 $x\in\mathbb R^3$ 上都为零.(相比之下,非零的非齐次多项式 $p(x):=x_1^2+x_2^2+x_3^2-1$ 在 $S^2$ 上恒等于零.)因此,我们可以自由地在把 $V_l$ 的元素看作 $S^2$ 上的函数和把它们看作 $\mathbb R^3$ 上的函数之间来回转换.

众所周知,Laplace 算子 $\Delta$ 与旋转交换.因此,每个 $V_l$ 都在旋转群的作用下不变.我们最终将看到,$V_l$ 在这个作用下是不可约的.

所有次数为 $0$ 或 $1$ 的齐次多项式都是调和的.因此,$V_0$ 由 $S^2$ 上的常值函数组成,而 $V_1$ 由函数 $x_1,x_2,x_3$ 在 $S^2$ 上的限制张成.另一方面,次数为 $2$ 的齐次多项式空间是六维的,而次数为 $2$ 的齐次调和多项式空间由下面五个多项式张成:
\[
 x_1x_2,\quad x_2x_3,\quad x_3x_1,\quad x_1^2-x_2^2,\quad x_2^2-x_3^2.
\]
多项式 $x_1^2-x_3^2$ 也是调和的,但它只是前两个差式 $x_1^2-x_2^2$ 与 $x_2^2-x_3^2$ 的和.

\begin{theorem}[$L^2(S^2)$ 中球谐函数的分解]\label{thm:spherical-harmonics-decomposition}
空间 $V_l$ 具有下列性质:

\begin{enumerate}
\item 每个 $V_l$ 的维数为 $2l+1$.
\item 每个 $V_l$ 在旋转群作用下不变,并且在该作用下不可约.
\item 当 $l\neq m$ 时,$V_l$ 与 $V_m$ 在 $L^2(S^2)$ 中正交.
\item Hilbert 空间 $L^2(S^2)$ 分解为所有 $V_l$ 的正交直和,其中 $l$ 遍历非负整数.
\end{enumerate}
\end{theorem}

本节余下部分将证明定理~\ref{thm:spherical-harmonics-decomposition}.我们将依次证明一系列引理以及这些引理的推论.

\begin{lemma}[多项式空间上的内积]\label{lem:polynomial-inner-product}
令 $\mathcal P$ 表示具有复系数的 $\mathbb R^3$ 上多项式空间.则 $\mathcal P$ 上存在内积 $\langle\cdot,\cdot\rangle_{\mathcal P}$,满足
\[
 \langle p,\Delta q\rangle_{\mathcal P}=\langle |x|^2p,q\rangle_{\mathcal P},
\]
其中 $|x|^2=x_1^2+x_2^2+x_3^2$.
\end{lemma}

\begin{proof}
虽然可以用组合方法构造所需内积,但也可以给出一个解析构造.$\mathbb R^3$ 上的每个多项式 $p$ 都有一个显然的全纯延拓到 $\mathbb C^3$,记为 $p_{\mathbb C}$.定义
\[
 \langle p,q\rangle_{\mathcal P}
 =\int_{\mathbb C^3}p_{\mathbb C}(z)\overline{q_{\mathbb C}(z)}\,\frac{e^{-|z|^2/2}}{\pi^{3/2}}\,d^6z,
\]
这正是把 $p_{\mathbb C}$ 与 $q_{\mathbb C}$ 看作 Segal--Bargmann 空间 $\mathcal H L^2(\mathbb C^3,\mu_1)$ 的元素时的内积.根据引理~\ref{lem:segal-gaussian-adjoint},对所有 $p,q\in\mathcal P$ 及 $j=1,2,3$,有
\[
 \int_{\mathbb C^3}p_{\mathbb C}(z)\frac{\partial q_{\mathbb C}}{\partial z_j}(z)\frac{e^{-|z|^2/2}}{\pi^{3/2}}\,d^6z
 =\int_{\mathbb C^3}z_jp_{\mathbb C}(z)q_{\mathbb C}(z)\frac{e^{-|z|^2/2}}{\pi^{3/2}}\,d^6z.
\]
这个关系意味着
\[
 \left\langle p,\frac{\partial q}{\partial x_j}\right\rangle_{\mathcal P}
 =\langle x_jp,q\rangle_{\mathcal P},
\]
于是所需的内积性质立即成立.
\end{proof}

标准的初等组合计算表明,满足 $l_1+l_2+l_3=l$ 的有序三元组 $(l_1,l_2,l_3)$ 的数量为 $(l+2)(l+1)/2$.由于满足 $l_1+l_2+l_3=l$ 的单项式 $x_1^{l_1}x_2^{l_2}x_3^{l_3}$ 构成 $\mathcal P_l$ 的一组基,故
\[
 \dim\mathcal P_l=\frac{(l+2)(l+1)}2,
\]
其中 $\mathcal P_l$ 表示次数为 $l$ 的齐次多项式空间.

\begin{corollary}[调和多项式空间的维数]\label{cor:dim-spherical-harmonics}
若 $\mathcal P_l$ 表示 $\mathbb R^3$ 上次数为 $l$ 的齐次多项式空间,则当 $l\ge2$ 时,Laplace 算子 $\Delta$ 将 $\mathcal P_l$ 映到 $\mathcal P_{l-2}$ 上且为满射.因此,对所有 $l\ge2$,有
\[
 \begin{aligned}
 \dim V_l&=\dim\mathcal P_l-\dim\mathcal P_{l-2}\\
 &=\frac{(l+2)(l+1)}2-\frac{l(l-1)}2=2l+1.
 \end{aligned}
\]
\end{corollary}

\begin{proof}
用引理~\ref{lem:polynomial-inner-product} 给出的内积装备有限维空间 $\mathcal P_l$ 与 $\mathcal P_{l-2}$.容易看出,对于从一个有限维内积空间到另一个有限维内积空间的线性映射,“像空间的正交补等于伴随映射的核”这一说法成立.应用于 $\Delta:\mathcal P_l\to\mathcal P_{l-2}$ 时,注意到 $\Delta$ 的伴随映射是乘以 $|x|^2$,而这个映射显然是单射,因为 $x_1^2+x_2^2+x_3^2$ 只在原点为零.因此,$\Delta$ 的像空间的正交补为 $\{0\}$.由于这些空间有限维,$\Delta$ 必为满射.
\end{proof}

\begin{corollary}[齐次多项式的调和分解]\label{cor:harmonic-decomposition-polynomial}
设 $l$ 为非负整数,并令 $k=l/2$(当 $l$ 为偶数时)或 $k=(l-1)/2$(当 $l$ 为奇数时).则每个 $p\in\mathcal P_l$ 都可以分解为
\[
 p(x)=p_0(x)+|x|^2p_1(x)+|x|^4p_2(x)+\cdots+|x|^{2k}p_k(x),
\]
其中每个 $p_j(x)$ 都是次数为 $l-2j$ 的齐次调和多项式.特别地,$p$ 在 $S^2$ 上的限制满足
\[
 p|_{S^2}=(p_0+p_1+\cdots+p_k)|_{S^2},
\]
这里 $p_0+p_1+\cdots+p_k$ 是一个非齐次调和多项式.
\end{corollary}

给定任意多项式 $p$(不必齐次),对 $p$ 的每个齐次部分应用推论~\ref{cor:harmonic-decomposition-polynomial},可知存在一个调和多项式 $\widetilde p$,使得 $p$ 与 $\widetilde p$ 在 $S^2$ 上具有相同的限制.

\begin{proof}
对 $l$ 归纳.若 $l=0$ 或 $l=1$,则每个 $p\in\mathcal P_l$ 都是调和的,所需分解就是 $p=p_0$.现在设 $l\ge2$,并假设所有小于 $l$ 的次数都已得到结论.由引理~\ref{lem:polynomial-inner-product},$\mathcal P_l$ 是 Laplace 算子的核与 $\mathcal P_{l-2}$ 在乘以 $|x|^2$ 后的像的正交直和.因此,每个 $p\in\mathcal P_l$ 都可以写成
\[
 p=p_0+|x|^2q_0,
\]
其中 $p_0$ 是调和的,$q_0\in\mathcal P_{l-2}$.由归纳假设,$q_0$ 具有所需形式的分解;将该分解代入 $p=p_0+|x|^2q_0$,便得到 $p$ 的所需分解.
\end{proof}

为了证明 $V_l$ 在 $SO(3)$ 的作用下不可约,我们转到 Lie 代数.正如已经指出的,齐次多项式到球面限制的映射是单射的,因此可以把 $V_j$ 的元素看作 $\mathbb R^3$ 上的多项式;此时,与 $\Pi$ 对应的 Lie 代数作用 $\pi$ 可用通常的角动量算子表示.

\begin{lemma}[球谐空间中的最高权向量]\label{lem:spherical-harmonics-highest-weight}
沿用定理~\ref{thm:so3-irreducible-standard-form} 的记号,令 $L_3=i\pi(F_3)=\widetilde J_3$,令
\[
 L_+=i\pi(F_1)-\pi(F_2)=\widetilde J_1+i\widetilde J_2.
\]
对任意非负整数 $l$,多项式
\[
 p(x_1,x_2,x_3):=(x_1+ix_2)^l
\]
属于 $V_l$,并满足
\[
 L_3p=lp,\qquad L_+p=0.
\]
\end{lemma}

\begin{proof}
因为 $p$ 与 $x_3$ 无关,并且作为 $z:=x_1+ix_2$ 的函数是全纯的,所以它自动是调和多项式;也可以直接计算验证这一点.另一方面,作用 $L_3$ 得
\[
\begin{aligned}
L_3p
&=-i\left(x_1\frac{\partial}{\partial x_2}-x_2\frac{\partial}{\partial x_1}\right)(x_1+ix_2)^l\\
&=-i\left[x_1\,l(x_1+ix_2)^{l-1}i-x_2\,l(x_1+ix_2)^{l-1}\right]\\
&=l(x_1+ix_2)^l.
\end{aligned}
\]
最后,应用 $L_+=i\pi(F_1)-\pi(F_2)$ 得
\[
\begin{aligned}
L_+p
&=\left[-ix_2\frac{\partial}{\partial x_3}-x_3\frac{\partial}{\partial x_2}
 +x_3\frac{\partial}{\partial x_1}-x_1\frac{\partial}{\partial x_3}\right]p\\
&=-i[-x_3l(x_1+ix_2)^{l-1}i]+x_3l(x_1+ix_2)^{l-1}\\
&=0,
\end{aligned}
\]
结论得证.
\end{proof}

\begin{corollary}[$V_l$ 的不可约性]\label{cor:spherical-harmonics-irreducible}
空间 $V_l$ 在 $SO(3)$ 的作用下不可约.
\end{corollary}

\begin{proof}
由命题~\ref{prop:so3-cyclic-chain},对上引理中的多项式 $p$ 反复应用 $L_-$,可以得到长度为 $2l+1$ 的特征向量“链”.这些特征向量张成一个维数为 $2l+1$ 的不可约不变子空间.由于已经证明 $\dim V_l=2l+1$,这条链中的向量必张成整个 $V_l$,从而 $V_l$ 不可约.
\end{proof}

现在证明本节的主要结果.

\begin{proof}[定理~\ref{thm:spherical-harmonics-decomposition} 的证明]
定理的第 1 点和第 2 点分别已经由推论~\ref{cor:dim-spherical-harmonics} 与推论~\ref{cor:spherical-harmonics-irreducible} 证明.现在,每个 $V_l$ 都是 $SO(3)$ 的不可约表示,而且由于它们的维数各不相同,任意两个 $V_l$ 都不可能同构.因此由第 16 章习题 19,$l\neq m$ 时,$V_l$ 与 $V_m$ 必须在 $L^2(S^2)$ 中正交,这就得到第 3 点.

最后,根据 Stone--Weierstrass 定理以及定理~\ref{app:thm:stone-complex} 的稠密性结论,$\mathbb R^3$ 上多项式限制到 $S^2$ 后得到的函数在 $L^2(S^2)$ 中构成稠密子空间.但推论~\ref{cor:harmonic-decomposition-polynomial} 表明,多项式在 $S^2$ 上的限制空间与调和多项式在 $S^2$ 上的限制空间相同.因此,各个 $V_j$ 的张成空间在 $L^2(S^2)$ 中稠密,得到第 4 点.
\end{proof}

\section{在 $L^2(\mathbb R^3)$ 中实现这些表示}
\label{sec:representations-in-l2-r3}

对于 $\mathbb R^3$ 上的齐次多项式,从 $\mathbb R^3$ 到 $S^2$ 的限制映射是单射的.因此,我们既可以把 $V_l$ 看作 $S^2$ 上的函数空间(如上一节),也可以把它看作 $\mathbb R^3$ 上的函数空间.本节中,我们将 $V_l$ 记为次数为 $l$ 的齐次调和多项式空间.

\begin{definition}[径向因子产生的表示空间]\label{def:radial-representation-subspace}
设 $l$ 是非负整数,$f$ 是 $(0,\infty)$ 上的可测函数,并满足
\[
 \int_0^\infty |f(r)|^2r^{2l+2}\,dr<\infty.
\label{eq:radial-factor-integrability}
\]
令 $V_{l,f}\subset L^2(\mathbb R^3)$ 为所有形如
\[
 \psi(x)=p(x)f(|x|),
\label{eq:radial-representation-function}
\]
其中 $p\in V_l$ 的函数所组成的空间.
\end{definition}

对 $f(r)$ 的这个条件,正好可以保证 $\psi(x)$ 属于 $L^2(\mathbb R^3)$(在球坐标中计算 $L^2$ 范数即可).

定义~\ref{def:radial-representation-subspace} 并不是物理学家通常使用的定义.在物理文献中,人们会考虑形如
\[
 \psi(x)=Y_{lm}(\theta,\phi)g(r)
\label{eq:spherical-harmonic-radial-form}
\]
的函数,其中 $r,\theta,\phi$ 是通常的球坐标,$Y_{lm}$ 是某个次数为 $l$ 的齐次调和多项式在球面上的限制,并用球坐标表示出来.(差一个归一化因子,$Y_{lm}$ 可由定理~\ref{thm:so3-irreducible-standard-form} 中的 $V_l$ 基得到.)因此,当我们沿 $\mathbb R^3$ 中的一条射线移动时,只有 $g(r)$ 的值发生变化.相比之下,在式~\eqref{eq:radial-representation-function} 中,因子 $p(x)$ 还会沿射线贡献一个 $r^l$ 因子.物理学家的表达式可以在直角坐标中写成
\[
 \psi(x)=Y_{lm}\left(\frac{x}{|x|}\right)g(|x|)
 =Y_{lm}(x)\frac{g(|x|)}{|x|^l}.
\label{eq:spherical-harmonic-rectangular-form}
\]
在计算上,式~\eqref{eq:radial-representation-function} 比式~\eqref{eq:spherical-harmonic-rectangular-form} 更方便;事实上,在分析氢原子时,物理学家会在计算后面的某个阶段乘上 $r^l$,正是为了让相关微分方程呈现更简单的形式.

\begin{proposition}[$L^2(\mathbb R^3)$ 中的有限维不可约表示]\label{prop:radial-representation-classification}
每个形如 $V_{l,f}\subset L^2(\mathbb R^3)$ 的空间在 $SO(3)$ 作用下不变且不可约.反过来,$L^2(\mathbb R^3)$ 中每个有限维、不可约且 $SO(3)$ 不变的子空间,都具有某个非负整数 $l$ 和满足式~\eqref{eq:radial-factor-integrability} 的函数 $f$ 所给出的 $V_{l,f}$ 形式.
\end{proposition}

\begin{proof}
由于因子 $f(|x|)$ 在旋转下不变,$SO(3)$ 的作用只影响函数 $p$.因此,作为 $SO(3)$ 的表示,$V_{l,f}$ 与空间 $V_l$ 同构,而后者由定理~\ref{thm:spherical-harmonics-decomposition} 知是不可约的.

对于反方向,$\mathbb R^3$ 上的 Lebesgue 测度分解为球面面积测度与 $(0,\infty)$ 上测度 $4\pi r^2dr$ 的乘积.因此,根据标准的测度论结果(命题~\ref{prop:l2-product-decomposition}),$L^2(\mathbb R^3)$ 可以自然地分解为 $L^2(S^2)$ 与 $L^2((0,\infty))$ 的 Hilbert 张量积;张量积中形如 $f\otimes g$ 的向量对应于 $L^2(\mathbb R^3)$ 中形如 $f(\theta,\phi)g(r)$ 的函数.由于定理~\ref{thm:spherical-harmonics-decomposition} 将 $L^2(S^2)$ 分解为 $l=0,1,2,\ldots$ 时各个 $V_l$ 的直和,我们可以选取 $L^2((0,\infty))$ 的一组标准正交基 $\{g_k\}$,将 $L^2(\mathbb R^3)$ 分解为形如
\[
 V_{l,k}:=V_l\otimes g_k
\]
的空间之和.

现在设 $V$ 是 $L^2(\mathbb R^3)$ 中任意有限维、不可约且 $SO(3)$ 不变的子空间.令 $\pi_{l,k}:L^2(\mathbb R^3)\to V_{l,k}$ 为正交投影,并令 $\rho_{l,k}$ 为其在 $V$ 上的限制.容易看出,$\rho_{l,k}$ 是 $SO(3)$ 作用的交织映射.因此,由于 $V$ 和 $V_{l,k}$ 都不可约,Schur 引理告诉我们,每个 $\rho_{l,k}$ 要么为零,要么为同构.此外,对于不同的 $l$,空间 $V_l$ 互不同构,所以不能同时有 $\rho_{l,k}$ 和 $\rho_{l',k'}$ 非零(其中 $l\neq l'$).另一方面,$\rho_{l,k}$ 不可能对所有 $k,l$ 都为零,因为各个 $V_{l,k}$ 的张成空间是整个 $L^2(\mathbb R^3)$.因此,存在某个 $l_0$,使得某个 $k_0$ 满足 $\rho_{l_0,k_0}\neq0$,同时对所有 $l\neq l_0$ 都有 $\rho_{l,k}=0$.

再次应用 Schur 引理,$\rho_{l_0,k}(\rho_{l_0,k_0})^{-1}$ 对每个 $k$ 都必须是某个常数 $c_k$ 乘以恒等映射.给定 $\psi\in V$,令 $v$ 是 $V$ 中唯一满足
\[
 \rho_{l_0,k_0}(\psi)=v\otimes g_{k_0}
\]
的元素.于是对每个 $k$,都有
\[
 \rho_{l_0,k}(\psi)=c_k(v\otimes g_k).
\]
又因为 $l\neq l_0$ 时 $\rho_{l,k}(\psi)=0$,故 $\psi$ 必须具有 $v\otimes g$ 的形式,其中
\[
 g=\sum_k c_kg_k.
\]
这对每个 $\psi\in V$ 都成立,而且常数 $c_k$ 与 $\psi$ 无关,因此 $V=V_{l_0}\otimes g$,这正是式~\eqref{eq:spherical-harmonic-radial-form} 的形式.于是 $V$ 具有命题所述的形式,其中
\[
 f(r)=\frac{g(r)}{r^{l_0}}.
\]
\end{proof}

还可以进一步证明,$L^2(\mathbb R^3)$ 中每个闭的、$SO(3)$ 不变的子空间都分解为有限维不可约、$SO(3)$ 不变子空间的正交直和.这个结果只是紧拓扑群的强连续酉表示的一般定理的一个特例(例如见文献 [10] 第 5 章).既然我们已经知道 $L^2(\mathbb R^3)$ 本身是有限维不可约不变子空间的直和,那么或许可以给出这个结果的初等证明,但我们不再沿这条路线展开.

\section{自旋}
\label{sec:spin}

我们按照自旋 $l$ 分类了 Lie 代数 $\mathfrak{so}(3)$ 的有限维不可约表示,其中 $l$ 是算子 $L_3=i\pi(F_3)$ 的最大特征值.$l$ 的可能取值是非负整数 $0,1,2,\ldots$ 以及正半整数 $1/2,3/2,\ldots$.然而,在 $L^2(S^2)$ 和 $L^2(\mathbb R^3)$ 中,我们只找到整数自旋的不可约表示.半整数自旋表示没有出现的原因很容易理解:它们不对应于 $SO(3)$ 的任何表示.由于 $L^2(S^2)$ 和 $L^2(\mathbb R^3)$ 都带有 $SO(3)$ 的自然酉作用 $\Pi$,任何在相应 Lie 代数表示 $\pi$ 下不变的有限维子空间,也会在 $\Pi$ 下不变,因而构成 $SO(3)$ 的表示.

虽然 $\mathfrak{so}(3)$ 的半整数表示不能积分为 $SO(3)$ 的表示,但它们可以像命题~\ref{prop:so3-group-representations} 的证明中那样,积分为 $SO(3)$ 的万能覆盖 $SU(2)$ 的表示.对半整数 $l$,相应的 $SU(2)$ 表示 $\Pi_l'$ 满足 $\Pi_l'(-I)=-I$,这意味着 $\Pi_l'$ 不能下降为
\[
 SO(3)\cong SU(2)/\{I,-I\}.
\]
不过,如果我们考虑射影表示,就会发现 $[-I]$ 是射影酉群 $PU(V)$ 中的单位元.因此,即使 $l$ 是半整数,我们仍然得到 $SO(3)$ 的一个良定义射影表示,满足
\[
 \Pi_l(e^{tX})=[e^{t\pi_l(X)}]
\]
对所有 $X\in\mathfrak{so}(3)$ 成立,其中 $[U]$ 表示 $U\in U(V)$ 在 $PU(V)$ 中的像.

通常认为,宇宙的物理规律在旋转群 $SO(3)$ 下不变.这并不意味着我们从不考虑没有旋转对称性的模型;例如,磁场中的氢原子所处的局部环境会破坏氢原子的旋转对称性.然而,如果把氢原子和磁场一起旋转,问题的物理规律并不会改变.在量子力学中,旋转对称性意味着:宇宙的 Hilbert 空间上应存在一个与 Hamilton 算子对易的 $SO(3)$ 射影酉表示.

现在,宇宙的 Hilbert 空间(如果确实存在这样的空间)由每一种粒子对应的 Hilbert 空间组合而成.因此,我们预期单个粒子的 Hilbert 空间也带有 $SO(3)$ 的射影酉表示.

单个粒子的 Hilbert 空间最简单的可能性是 $L^2(\mathbb R^3)$,它当然带有我们一直讨论的 $SO(3)$ 的普通酉作用.然而,根据各种实验观察,物理学家提出了对单个粒子 Hilbert 空间的修正,将“内部自由度”纳入其中.该设想是:对于每一种粒子,量子 Hilbert 空间应具有
\[
 L^2(\mathbb R^3)\widehat\otimes V
\]
的形式,其中 $V$ 是一个有限维 Hilbert 空间,并承载 $SO(3)$ 的不可约射影酉表示.这里 $\widehat\otimes$ 是 Hilbert 张量积(见附录~\ref{app:subsec:hilbert-tensor}).$SO(3)$ 在 $V$ 上的射影作用描述旋转群对粒子内部自由度的作用.

根据命题~\ref{prop:projective-representation-lie-lift},空间 $V$ 承载一个迹为零的普通 Lie 代数表示 $\pi:\mathfrak{so}(3)\to\mathfrak{gl}(V)$.按物理学的通常术语,算子 $L_3:=i\pi(F_3)$ 在 $V$ 上的最大特征值 $l$ 称为粒子的自旋.于是用 $V_l$ 表示自旋为 $l$ 的空间.例如,电子是“自旋 $1/2$”粒子,意味着单个电子的 Hilbert 空间为
\[
 L^2(\mathbb R^3)\widehat\otimes V_{1/2},
\]
其中 $V_{1/2}$ 是二维的 $SO(3)$ 射影表示.

容易看出,给定群的两个射影酉表示的张量积仍是该群的射影酉表示.(相反,两个射影酉表示的直和一般不再是射影酉表示.)在当前情形中,可以把 $L^2(\mathbb R^3)$ 看作 $SU(2)$ 的酉表示 $\Pi$,它经由 $SO(3)$ 分解,也就是满足 $\Pi(-I)=I$.另一方面,可以把 $V_l$ 看作 $SU(2)$ 的酉表示 $\Pi_l$,其中
\[
 \Pi_l(-I)=\begin{cases}I,&l\text{ 为整数},\\-I,&l\text{ 为半整数}.\end{cases}
\]
于是 $L^2(\mathbb R^3)\widehat\otimes V_l$ 带有 $SU(2)$ 的表示 $\Pi\otimes\Pi_l$,并且
\[
 (\Pi\otimes\Pi_l)(-I)=\pm I.
\]
因此从射影的意义看,$\Pi\otimes\Pi_l$ 可以下降为 $SO(3)$ 的表示.

\begin{note}[自旋]\label{note:spin-summary}
每一种粒子都有一个“自旋”$l$,其中 $l$ 是非负整数或半整数.这类粒子的 Hilbert 空间为
\[
 L^2(\mathbb R^3)\widehat\otimes V_l,
\]
其中 $V_l$ 是 $SO(3)$ 的维数为 $2l+1$ 的不可约射影表示.
\end{note}

由于 $V_l$ 是有限维的,Hilbert 张量积 $L^2(\mathbb R^3)\widehat\otimes V_l$ 与 $L^2(\mathbb R^3)$ 和 $V_l$ 的代数张量积相同.

\begin{definition}[玻色子与费米子]\label{def:boson-fermion}
自旋为整数的粒子称为玻色子,自旋为半整数的粒子称为费米子.
\end{definition}

要理解整数自旋与半整数自旋之间区别的重要性,需要考察描述多个同类粒子的 Hilbert 空间,例如五个电子的 Hilbert 空间.第 19 章将讨论这个主题.

\section{表示的张量积:“角动量相加”}
\label{sec:tensor-products-angular-momentum}

设 $V_l$ 与 $V_m$ 分别是维数为 $2l+1$ 和 $2m+1$ 的 $\mathfrak{so}(3)$ 不可约表示.如第 16 章~\ref{sec:new-representations-from-old} 节所述,张量积空间 $V_l\otimes V_m$ 可以看作 $\mathfrak{so}(3)$ 的另一个表示.除非 $l$ 与 $m$ 中至少有一个为零,否则 $V_l\otimes V_m$ 不可约.因此,我们关心把 $V_l\otimes V_m$ 分解为不可约不变子空间的直和.在 $V_l$ 是 $L^2(\mathbb R^3)$ 中的不可约 $SO(3)$ 不变子空间、而 $V_m$ 是粒子内部自由度空间的情形,这个分解将帮助我们把带自旋粒子的 Hilbert 空间分解为不可约的 $SO(3)$ 不变子空间.

\begin{proposition}[$V_l\otimes V_{1/2}$ 的分解]\label{prop:so3-tensor-half}
设 $V_{1/2}$ 是维数为 $2$ 的 $\mathfrak{so}(3)$ 不可约表示,$V_l$ 是维数为 $2l+1$ 的 $\mathfrak{so}(3)$ 不可约表示,其中 $l$ 为非负整数或半整数.若 $l=0$,则 $V_l\otimes V_{1/2}$ 不可约;若 $l>0$,则
\[
 V_l\otimes V_{1/2}\cong V_{l+1/2}\oplus V_{l-1/2}.
\]
\end{proposition}

\begin{proof}
若 $l=0$,容易看出 $V_l\otimes V_{1/2}$ 同构于不可约的 $V_{1/2}$.以下设 $l>0$.

令 $L_+,L_-,L_3$ 为利用表示 $\pi_l$ 按定理~\ref{thm:so3-irreducible-standard-form} 构造的算子,令 $\sigma_+,\sigma_-,\sigma_3$ 为利用表示 $\pi_{1/2}$ 构造的相应算子.按照第 16 章~\ref{sec:new-representations-from-old} 节,在 $V_l\otimes V_{1/2}$ 上定义
\[
 J_+=L_+\otimes I+I\otimes\sigma_+,\qquad
 J_-=L_-\otimes I+I\otimes\sigma_-,\qquad
 J_3=L_3\otimes I+I\otimes\sigma_3.
\label{eq:angular-momentum-tensor-product-operators}
\]
令 $\{v_0,\ldots,v_{2l}\}$ 是 $V_l$ 的一组定理~\ref{thm:so3-irreducible-standard-form} 所述基,令 $\{e_0,e_1\}$ 是 $V_{1/2}$ 的类似基.$v_j\otimes e_k$ 构成 $V_l\otimes V_{1/2}$ 的一组基.$J_3$ 的特征值为
\[
 (l-j)+\left(\frac12-k\right),
\qquad j=0,1,\ldots,2l,\quad k=0,1.
\]
因此,$J_3$ 的特征值从 $l+1/2$ 取到 $-(l+1/2)$.其中 $l+1/2$ 与 $-(l+1/2)$ 各只出现一次;其他特征值 $\lambda$ 都出现两次,分别写成 $(\lambda-1/2)+1/2$ 与 $(\lambda+1/2)-1/2$.

向量 $v_0\otimes e_0$ 是 $J_3$ 的最大特征值 $l+1/2$ 对应的特征向量,因此 $J_+(v_0\otimes e_0)=0$.由命题~\ref{prop:so3-cyclic-chain},反复应用 $J_-$ 会得到长度为 $2l+2$ 的特征向量链,这些向量张成一个同构于 $V_{l+1/2}$ 的不可约不变子空间 $W_0$.

另一方面,由命题~\ref{prop:so3-unitary-inner-product},可以在 $V_l$ 和 $V_{1/2}$ 上分别选取内积,使得 $\pi_l$ 与 $\pi_{1/2}$ 成为“酉”的,即对所有 $X\in\mathfrak{so}(3)$ 都有 $\pi(X)^*=-\pi(X)$.在 $V_l\otimes V_{1/2}$ 上使用由这两个内积按附录~\ref{app:subsec:hilbert-tensor} 诱导的自然内积,则 $\pi_l\otimes\pi_{1/2}$ 也是酉的.因此,不变子空间 $W_0$ 的正交补也不变.由于除了最大和最小特征值外,$J_3$ 的所有特征值重数都是 $2$,所以 $J_3$ 在 $W_0^\perp$ 中的最大特征值为 $l-1/2$.取 $w_0\in W_0^\perp$ 为对应于特征值 $l-1/2$ 的特征向量.反复应用降算子
\[
 J_-=L_-\otimes I+I\otimes\sigma_-
\]
会得到长度为 $2l$ 的特征向量链.这些特征向量张成 $V_l\otimes V_{1/2}$ 的一个维数为 $2l$ 的不可约不变子空间 $W_1$.由于
\[
 \dim W_0+\dim W_1=4l+2=\dim(V_l\otimes V_{1/2}),
\]
必有 $W_1=W_0^\perp$,证明完成.
\end{proof}

由于电子是“自旋 $1/2$”粒子,根据第~\ref{sec:spin} 节,单个电子的 Hilbert 空间为 $L^2(\mathbb R^3)\widehat\otimes V_{1/2}$,其中 $V_{1/2}$ 是维数为 $2$ 的不可约 $SO(3)$ 射影酉表示.另一方面,第~\ref{sec:representations-in-l2-r3} 节说明,对于 $l=0,1,2,\ldots$,可以找到 $L^2(\mathbb R^3)$ 中维数为 $2l+1$ 的不可约 $SO(3)$ 不变子空间 $V_{l,f}$,其中 $f$ 是任意径向函数.将命题~\ref{prop:so3-tensor-half} 应用于 $V_l=V_{l,f}$,就得到带自旋粒子的 Hilbert 空间中的不可约 $SO(3)$ 不变子空间.比如,在分析氢原子的精细结构时,找到这些子空间是必不可少的.

当 $V_l$ 是 $L^2(\mathbb R^3)$ 中的 $SO(3)$ 不变子空间时,式~\eqref{eq:angular-momentum-tensor-product-operators} 中的 $J_3$ 在物理文献中写作
\[
 J_3=L_3+\sigma_3,
\label{eq:total-angular-momentum-third-component}
\]
其中默认 $L_3$ 作用于张量积的第一因子,$\sigma_3$ 作用于第二因子(也就是说,与恒等算子的张量积被省略不写).这里 $L_3=L_3\otimes I$ 是通常的轨道角动量算子,而 $\sigma_3=I\otimes\sigma_3$ 描述基元素 $F_3\in\mathfrak{so}(3)$ 在 $V_{1/2}$ 上的作用.类似地,$L_\pm$ 与 $\sigma_\pm$ 也有相应含义.物理学家把对 $\mathfrak{so}(3)$ 表示的张量积分析称为“角动量相加”,正是因为有式~\eqref{eq:total-angular-momentum-third-component} 这样的公式.在这个语境下,算子 $L_3$ 称为轨道角动量算子,算子 $\sigma_3$ 称为自旋角动量算子.

现在记录 $\mathfrak{so}(3)$ 不可约表示张量积的一般结果.

\begin{proposition}[不可约 $\mathfrak{so}(3)$ 表示的张量积分解]\label{prop:so3-tensor-product-decomposition}
对任意 $j=0,1/2,1,\ldots$,令 $V_j$ 表示唯一的维数为 $2j+1$ 的 $\mathfrak{so}(3)$ 不可约表示.若 $l\ge m$,则
\[
 V_l\otimes V_m\cong V_{l+m}\oplus V_{l+m-1}\oplus\cdots\oplus V_{l-m+1}\oplus V_{l-m}.
\label{eq:so3-tensor-product-decomposition}
\]
\end{proposition}

这个结果的证明与命题~\ref{prop:so3-tensor-half} 的证明类似,此处略去;见文献 [21] 附录 D 的定理 D.1.分解~\eqref{eq:so3-tensor-product-decomposition} 的一个重要性质是:右端出现的每个不可约表示都只出现一次.这个性质正是 Wigner--Eckart 定理证明中的关键思想.详细内容见文献 [21] 附录 D.

\section{向量与向量算子}
\label{sec:vectors-vector-operators}

\begin{definition}[按向量变换的函数]\label{def:classical-vector-transform}
若函数 $c:\mathbb R^3\times\mathbb R^3\to\mathbb R^3$ 满足对所有 $R\in SO(3)$,
\[
 c(Rx,Rp)=R(c(x,p)),
\label{eq:classical-vector-transformation}
\]
则称 $c$ 按向量变换.
\end{definition}

在物理文献中,有时用“是一个向量”来代替“按向量变换”.

注意,定义~\ref{def:classical-vector-transform} 只考虑 $c$ 在 $SO(3)$ 元素下的变换性质,而不是一般 $O(3)$ 元素下的变换性质.如果 $c$ 对所有 $R\in O(3)$ 都满足式~\eqref{eq:classical-vector-transformation}(而不仅是 $SO(3)$ 中的 $R$),则称 $c$ 为“真向量”;如果对 $R\in O(3)\setminus SO(3)$ 满足 $c(Rx,Rp)=-R(c(x,p))$,则称 $c$ 为“赝向量”.对本书的目的而言,没有必要区分真向量和赝向量.

位置函数 $c_1(x,p):=x$、动量函数 $c_2(x,p):=p$,以及角动量函数 $c_3(x,p):=x\times p$,都是按向量变换的函数.(叉积在旋转下变换的性质是标准事实.)另一个典型例子是
\[
 c(x,p)=(x\cdot p)|x|^{-1}(x\times p).
\]

\begin{proposition}[经典向量的 Poisson 括号关系]\label{prop:classical-vector-poisson-relations}
令 $j(x,p)=x\times p$ 为 $\mathbb R^3\times\mathbb R^3$ 上的角动量函数.若光滑函数 $c:\mathbb R^3\times\mathbb R^3\to\mathbb R^3$ 按向量变换,则对 $k=1,2,3$,有
\[
 \{c_k,j_k\}=0.
\label{eq:classical-vector-poisson-diagonal}
\]
此外,
\[
 \{c_1,j_2\}=\{j_1,c_2\}=c_3,
\label{eq:classical-vector-poisson-cyclic}
\]
以及通过对指标作循环置换得到的其他关系.
\end{proposition}

\begin{proof}
令 $R(\theta)$ 表示 $(x_1,x_2)$ 平面内逆时针旋转角 $\theta$ 的旋转.将式~\eqref{eq:classical-vector-transformation} 用于 $R=R(\theta)$,并只看向量的第一分量,得到
\[
 c_1(R(\theta)x,R(\theta)p)=c_1(x,p)\cos\theta-c_2(x,p)\sin\theta.
\label{eq:classical-vector-rotation-component}
\]
正如命题~\ref{pro:angular-momentum-flow} 的证明中所述,Poisson 括号 $\{c_1,j_3\}$ 正好是式~\eqref{eq:classical-vector-rotation-component} 左端关于 $\theta$ 的导数在 $\theta=0$ 处的值.因此
\[
 \{c_1,j_3\}=-c_2,
\]
从而 $\{j_3,c_1\}=c_2$,这是对式~\eqref{eq:classical-vector-poisson-cyclic} 循环置换指标得到的关系之一.

另一方面,仍取 $R=R(\theta)$,但现在考察第三分量,有
\[
 c_3(R(\theta)x,R(\theta)p)=c_3(x,p).
\]
对 $\theta$ 求导并令 $\theta=0$,得到 $\{c_3,j_3\}=0$.其他 Poisson 括号关系同理计算.
\end{proof}

现在转向按向量变换的函数在量子力学中的对应物.

\begin{definition}[向量算子]\label{def:vector-operator}
对 $L^2(\mathbb R^3)$ 上的任意算子三元组 $C:=(C_1,C_2,C_3)$ 以及任意向量 $v\in\mathbb R^3$,定义
\[
 v\cdot C=\sum_{j=1}^3v_jC_j.
\label{eq:vector-operator-component}
\]
若算子三元组 $C$ 满足对所有 $R\in SO(3)$,
\[
 (Rv)\cdot C=\Pi(R)(v\cdot C)\Pi(R)^{-1},
\label{eq:vector-operator-transformation}
\]
则称 $C$ 为向量算子.
\end{definition}

这里 $\Pi$ 是定义~\ref{def:rotation-representation} 中 $SO(3)$ 在 $L^2(\mathbb R^3)$ 上的自然酉作用.我们以角动量为例,尝试理解这个定义的含义;角动量正是一个向量算子(如我们将看到的).算子 $\widehat J_1,\widehat J_2,\widehat J_3$ 分别表示 $\widehat J$ 在 $e_1,e_2,e_3$ 方向上的分量.更一般地,可以考察任意单位向量 $v$ 方向上的 $\widehat J$ 分量,它就是式~\eqref{eq:vector-operator-component} 定义的 $v\cdot\widehat J$.由于空间中没有优先方向,我们预期任意两个单位向量 $v_1,v_2$ 对应的算子 $v_1\cdot\widehat J$ 与 $v_2\cdot\widehat J$ 应当“相差一个旋转后是同一个算子”.具体而言,若旋转 $R$ 满足 $Rv_1=v_2$,则 $v_1\cdot\widehat J$ 与 $v_2\cdot\widehat J$ 应只相差 Hilbert 空间中 $R$ 的作用,即
\[
 v_2\cdot\widehat J=\Pi(R)(v_1\cdot\widehat J)\Pi(R)^{-1},
\]
这正是式~\eqref{eq:vector-operator-transformation} 在 $v=v_1,C=\widehat J$ 时的内容.

我们不讨论式~\eqref{eq:vector-operator-transformation} 是否仍对 $R\in O(3)\setminus SO(3)$ 成立.位置算子 $X$ 和动量算子 $P$ 很容易看出是向量算子.与经典情形一样,两个向量算子的叉积仍是向量算子(见第 18 章习题 7).特别地,角动量 $\widehat J=X\times P$ 是一个向量算子.

如果 $C_1,C_2,C_3$ 是无界算子,则定义~\ref{def:vector-operator} 还需要说明算子的定义域.最简单的处理方式是找一个 $L^2(\mathbb R^3)$ 的稠密子空间 $\mathcal V$,它包含于每个 $C_j$ 的定义域,并且在旋转下不变.此时,式~\eqref{eq:vector-operator-transformation} 被理解为应用于 $\mathcal V$ 中向量时成立.在许多情形下可以取 $\mathcal V=\mathcal S(\mathbb R^3)$.下面命题中,$\mathcal V$ 应满足某些技术性的定义域条件,以便对作用于 $\mathcal V$ 中向量的式~\eqref{eq:vector-operator-component} 求导.我们不展开这些条件的细节.

\begin{proposition}[向量算子的交换关系]\label{prop:vector-operator-commutators}
若 $C$ 是向量算子,则其分量满足
\[
 \frac1{i\hbar}[C_j,\widehat J_j]=0,\qquad j=1,2,3.
\label{eq:vector-operator-diagonal-commutator}
\]
此外,
\[
 \frac1{i\hbar}[C_1,\widehat J_2]=\frac1{i\hbar}[\widehat J_1,C_2]=C_3,
\label{eq:vector-operator-cyclic-commutator}
\]
以及对指标作循环置换得到的其他关系.
\end{proposition}

\begin{proof}
如命题~\ref{prop:classical-vector-poisson-relations} 的证明,令 $R(\theta)$ 表示 $(x_1,x_2)$ 平面内的旋转,令 $e_1=(1,0,0)$.在式~\eqref{eq:vector-operator-transformation} 中取 $R=R(\theta)$、$v=e_1$,得到
\[
 \Pi(R(\theta))C_1\Pi(R(\theta))^{-1}=C_1\cos\theta+C_2\sin\theta.
\label{eq:vector-operator-rotation-conjugation}
\]
然而 $R(\theta)=e^{\theta F_3}$,其中 $\{F_j\}$ 是第~\ref{sec:lie-algebra-of-matrix-group} 节所描述的 $\mathfrak{so}(3)$ 基.因此,对式~\eqref{eq:vector-operator-rotation-conjugation} 关于 $\theta$ 求导并令 $\theta=0$,得到
\[
 \pi(F_3)C_1-C_1\pi(F_3)=C_2.
\]
由于命题~\ref{prop:angular-momentum-lie-action} 给出 $\widehat J_3=i\hbar\pi(F_3)$,于是
\[
 \frac1{i\hbar}[\widehat J_3,C_1]=C_2,
\]
这正是对式~\eqref{eq:vector-operator-cyclic-commutator} 循环置换变量得到的关系之一.

另一方面,在式~\eqref{eq:vector-operator-transformation} 中取 $v=e_3$,得到
\[
 \Pi(R(\theta))C_3\Pi(R(\theta))^{-1}=C_3.
\]
对 $\theta$ 求导并令 $\theta=0$,得 $[\pi(F_3),C_3]=0$.所有其他关系都可以同样得到.

关于向量算子的更多信息,包括 Wigner--Eckart 定理,见文献 [21] 附录 D;另见习题 7.
\end{proof}

\section{练习}
\label{sec:chapter17-exercises}

\begin{exercise}[验证旋转向量场]\label{ex:rotation-vector-field}
验证式~\eqref{eq:j3-rotation-derivative} 中的向量场 $x_1\,\partial/\partial x_2-x_2\,\partial/\partial x_1$ 的表达式.
\end{exercise}

\begin{exercise}[验证升算子公式]\label{ex:so3-raising-induction}
在式~\eqref{eq:so3-raising-formula} 的证明中,利用归纳法和交换关系~\eqref{eq:so3-L-commutators} 验证式~\eqref{eq:so3-raising-formula}.
\end{exercise}

\begin{exercise}[证明 Casimir 算子命题]\label{ex:so3-casimir}
本题给出命题~\ref{prop:so3-casimir} 的一种证明.设 $(\pi,V_l)$ 是维数为 $2l+1$ 的 $\mathfrak{so}(3)$ 不可约表示,令 $C_\pi$ 为该命题中定义的 Casimir 算子.
\begin{enumerate}
\item 证明对 $j=1,2,3$,有 $[\pi(F_j),C_\pi]=0$.
\item 使用 Schur 引理证明存在某个 $\lambda\in\mathbb C$,使对所有 $v\in V$,$C_\pi v=\lambda v$.
\item 证明
\[
 C_\pi=-(L_3)^2+L_-L_++L_3,
\]
其中 $L_+,L_-,L_3$ 如定理~\ref{thm:so3-irreducible-standard-form} 所定义.
\item 通过在 $V$ 中适当选择的向量上计算 $C_\pi$,证明第 2 部分中的常数 $\lambda$ 等于 $-l(l+1)$.
\end{enumerate}
\end{exercise}

\begin{exercise}[构造 $\mathfrak{so}(3)$ 的不可约表示]\label{ex:so3-representation-commutators}
令 $l$ 为任意非负整数或半整数.规定向量 $\{v_0,v_1,\ldots,v_{2l}\}$ 构成向量空间 $V$ 的一组基,并用式~\eqref{eq:so3-irreducible-action} 定义 $V$ 上算子 $L_+,L_-,L_3$ 的作用.证明这些算子满足交换关系~\eqref{eq:so3-L-commutators}.

提示:在计算 $L_-$ 时,单独处理基向量 $v_{2l}$;在计算 $L_+$ 时,单独处理基向量 $v_0$.
\end{exercise}

\begin{exercise}[直和表示中的循环向量]\label{ex:so3-cyclic-vector-direct-sum}
设 $(\pi,V)$ 是维数为 $2$ 的 $\mathfrak{so}(3)$ 不可约表示,并取定理~\ref{thm:so3-irreducible-standard-form} 中的基 $\{v_0,v_1\}$.将 $V\oplus V$ 看作第 16 章~\ref{sec:new-representations-from-old} 节所述的 $\mathfrak{so}(3)$ 表示.令 $v=(v_0,v_1)$.证明包含 $v$ 的最小不变子空间是 $V\oplus V$.

\begin{note}
这说明 $V\oplus V$ 虽然不可约性失败,却仍然具有循环向量.
\end{note}
\end{exercise}

\begin{exercise}[计算 $V_1\otimes V_{1/2}$ 的基]\label{ex:so3-tensor-half-bases}
计算两个不可约不变子空间 $W_0\cong V_{3/2}$ 与 $W_0^\perp\cong V_{1/2}$ 的显式基.$W_0$ 或 $W_0^\perp$ 的每个基元素都应表示为命题~\ref{prop:so3-tensor-half} 的证明中向量 $v_j\otimes e_k$ 的线性组合.
\end{exercise}

\begin{exercise}[交织映射的唯一性]\label{ex:so3-intertwiner-uniqueness}
设 $V_l,V_m,V_n$ 分别是维数为 $2l+1,2m+1,2n+1$ 的 $\mathfrak{so}(3)$ 不可约表示.假设 $\Phi$ 与 $\Psi$ 是从 $V_l$ 到 $V_m\otimes V_n$ 的非零交织映射.证明存在 $c\in\mathbb C$,使 $\Phi=c\Psi$.

提示:使用命题~\ref{prop:so3-tensor-product-decomposition} 与 Schur 引理.

\begin{note}
这个结果与“不可约张量算子”的 Wigner--Eckart 定理密切相关.
\end{note}
\end{exercise}
